% -----------------------------
% PUCON User Manual (LaTeX)
% File: PUCON_User_Manual.tex
% Two versions included below:
%  1) Commented version (contains inline comments explaining each part)
%  2) Clean/uncommented version (ready to compile)
% -----------------------------

% -----------------------------
% 1) Commented version
% -----------------------------
% The commented version explains packages and structure. Remove the commented
% blocks or use the clean version below if you only need the source to compile.

% -----------------------------
% BEGIN COMMENTED VERSION
% -----------------------------
% Document class: article is simple and portable. For a longer manual consider
% book or report.
\documentclass[11pt,a4paper]{article}

% Basic packages
% geometry controls margins, graphicx loads images, hyperref for links,
% fancyhdr for headers/footers, caption for figure captions, setspace for
% spacing and enumitem for nicer lists.
\usepackage[margin=1in]{geometry} % adjust margins
\usepackage{graphicx}
\usepackage{hyperref}
\usepackage{fancyhdr}
\usepackage{caption}
\usepackage{setspace}
\usepackage{enumitem}
\usepackage{longtable}
\usepackage{titlesec}
\usepackage{tocloft}

% Header/Footer setup
\pagestyle{fancy}
\fancyhf{}
\lhead{PUCON}
\rhead{User Manual}
\cfoot{Page \thepage}

% Title formatting
\titleformat{\section}{\large\bfseries}{\thesection.}{0.5em}{}
\titleformat{\subsection}{\normalsize\bfseries}{\thesubsection.}{0.5em}{}

% Hyperlink setup
\hypersetup{
    colorlinks=true,
    linkcolor=blue,
    urlcolor=blue,
    pdftitle={PUCON User Manual},
    pdfauthor={PUCON Team}
}

% spacing
\onehalfspacing

% Begin document
\begin{document}

% Title page
\begin{titlepage}
    \centering
    {\Huge \bfseries PUCON User Manual \\[0.5em]}
    {\Large PUPCON Accreditation System \\}
    \\[2em]
    {\large Version: 1.0 \\
    Prepared by: PUCON Team \\
    Date: \today}
    \\
    \vfill
    % Optional logo (uncomment and replace path if needed)
    % 
\includegraphics[width=0.3\textwidth]{assets/logo.png}
\end{titlepage}

\tableofcontents
\newpage

% -----------------------------
% Sections: follow the structure provided earlier
% -----------------------------

\section{Introduction}
% Short description of the system and purpose of manual
PUCON (PUPCON Accreditation System) is a web-based platform designed to
streamline the accreditation process for coordinators, area chairs, and
evaluators. This manual explains how to access and use key features.

\section{System Requirements}
\begin{itemize}
    \item \textbf{Supported Browsers:} Google Chrome, Mozilla Firefox, Microsoft Edge
    \item \textbf{Internet:} Stable connection recommended (\textgreater= 5 Mbps)
    \item \textbf{Device:} Desktop or Laptop (recommended resolution: 1366x768 or higher)
\end{itemize}

\section{Accessing the System}
\begin{enumerate}
    \item Open your web browser.
    \item Navigate to: \url{https://pucon.example.com}.
    \item Enter your registered \textbf{Username} and \textbf{Password}.
    \item Click \textbf{Login}.
\end{enumerate}

\subsection{Forgot Password}
If you forgot your password, click the \textbf{Forgot Password} link and follow the
recovery instructions sent to your registered email.

\section{User Roles}
\begin{description}[leftmargin=!,labelwidth=3cm]
    \item[Accreditation Coordinator] Oversees and manages the entire accreditation process.
    \item[Area Chair] Uploads, reviews, and monitors documents for specific areas.
    \item[Evaluator] Reviews submissions and provides feedback.
    \item[Admin] Manages users and system settings.
\end{description}

\section{Main Features}
\subsection{Dashboard}
Displays pending tasks, recent activities, and system notifications.

\subsection{Document Management}
Upload, view, and download accreditation-related files. Organize documents by accreditation area.

\subsection{Evaluation Module}
Submit evaluation forms and view consolidated feedback and results.

\subsection{Reports}
Generate and download summary reports. Filter by date, area, or evaluator.

\section{Common Tasks}
\subsection{Uploading a Document}
\begin{enumerate}
    \item Go to \textbf{Document Management}.
    \item Select the appropriate \textbf{Accreditation Area}.
    \item Click \textbf{Upload File}.
    \item Browse and select your file.
    \item Click \textbf{Submit}.
\end{enumerate}

\subsection{Viewing Evaluation Results}
\begin{enumerate}
    \item Navigate to \textbf{Evaluation Module}.
    \item Select your submitted documents.
    \item Click \textbf{View Results} to see feedback and scores.
\end{enumerate}

\section{Troubleshooting}
\begin{itemize}
    \item \textbf{Login Issues:} Check your username/password. Use \textbf{Forgot Password} if needed.
    \item \textbf{File Upload Fails:} Ensure the file is within allowed size and format.
    \item \textbf{Page Not Loading:} Refresh the page or check your internet connection.
\end{itemize}

\section{Support}
For assistance, contact the PUCON Support Team at:
\begin{itemize}
    \item Email: \href{mailto:support@pucon.example.com}{support@pucon.example.com}
    \item Phone: +63-XXX-XXX-XXXX
\end{itemize}

% Optional: Appendix with screenshots or diagrams
\newpage
\appendix
\section{Appendix A: Screenshots}
% Example of including an inline image. Replace path as needed.
\begin{center}
    % Uncomment and set correct path to include diagram
    % \includegraphics[width=0.8\textwidth]{assets/architecture_dia_pupsched.png}
    \captionof{figure}{Architecture diagram (optional).}
\end{center}

\section{Appendix B: Change Log}
\begin{longtable}{p{3cm} p{10cm}}
\textbf{Version} & \textbf{Changes} \\
1.0 & Initial release of the user manual. \\
\end{longtable}

\end{document}

% -----------------------------
% END COMMENTED VERSION
% -----------------------------


% -----------------------------
% 2) Clean / Uncommented version (ready to compile)
% -----------------------------

% Use the content below if you prefer a compact file without explanatory comments.

% ----- BEGIN CLEAN VERSION -----

\documentclass[11pt,a4paper]{article}
\usepackage[margin=1in]{geometry}
\usepackage{graphicx}
\usepackage{hyperref}
\usepackage{fancyhdr}
\usepackage{caption}
\usepackage{setspace}
\usepackage{enumitem}
\usepackage{longtable}
\usepackage{titlesec}
\usepackage{tocloft}

\pagestyle{fancy}
\fancyhf{}
\lhead{PUCON}
\rhead{User Manual}
\cfoot{Page \thepage}

\titleformat{\section}{\large\bfseries}{\thesection.}{0.5em}{}
\titleformat{\subsection}{\normalsize\bfseries}{\thesubsection.}{0.5em}{}

\hypersetup{
    colorlinks=true,
    linkcolor=blue,
    urlcolor=blue,
    pdftitle={PUCON User Manual},
    pdfauthor={PUCON Team}
}

\onehalfspacing

\begin{document}

\begin{titlepage}
    \centering
    {\Huge \bfseries PUCON User Manual \\[0.5em]}
    {\Large PUPCON Accreditation System \\}
    \\[2em]
    {\large Version: 1.0 \\
    Prepared by: PUCON Team \\
    Date: \today}
    \\[2em]
\end{titlepage}

\tableofcontents
\newpage

\section{Introduction}
PUCON (PUPCON Accreditation System) is a web-based platform designed to
streamline the accreditation process for coordinators, area chairs, and
evaluators. This manual explains how to access and use key features.

\section{System Requirements}
\begin{itemize}
    \item \textbf{Supported Browsers:} Google Chrome, Mozilla Firefox, Microsoft Edge
    \item \textbf{Internet:} Stable connection recommended (\textgreater= 5 Mbps)
    \item \textbf{Device:} Desktop or Laptop (recommended resolution: 1366x768 or higher)
\end{itemize}

\section{Accessing the System}
\begin{enumerate}
    \item Open your web browser.
    \item Navigate to: \url{https://pucon.example.com}.
    \item Enter your registered \textbf{Username} and \textbf{Password}.
    \item Click \textbf{Login}.
\end{enumerate}

\subsection{Forgot Password}
If you forgot your password, click the \textbf{Forgot Password} link and follow the
recovery instructions sent to your registered email.

\section{User Roles}
\begin{description}[leftmargin=!,labelwidth=3cm]
    \item[Accreditation Coordinator] Oversees and manages the entire accreditation process.
    \item[Area Chair] Uploads, reviews, and monitors documents for specific areas.
    \item[Evaluator] Reviews submissions and provides feedback.
    \item[Admin] Manages users and system settings.
\end{description}

\section{Main Features}
\subsection{Dashboard}
Displays pending tasks, recent activities, and system notifications.

\subsection{Document Management}
Upload, view, and download accreditation-related files. Organize documents by accreditation area.

\subsection{Evaluation Module}
Submit evaluation forms and view consolidated feedback and results.

\subsection{Reports}
Generate and download summary reports. Filter by date, area, or evaluator.

\section{Common Tasks}
\subsection{Uploading a Document}
\begin{enumerate}
    \item Go to \textbf{Document Management}.
    \item Select the appropriate \textbf{Accreditation Area}.
    \item Click \textbf{Upload File}.
    \item Browse and select your file.
    \item Click \textbf{Submit}.
\end{enumerate}

\subsection{Viewing Evaluation Results}
\begin{enumerate}
    \item Navigate to \textbf{Evaluation Module}.
    \item Select your submitted documents.
    \item Click \textbf{View Results} to see feedback and scores.
\end{enumerate}

\section{Troubleshooting}
\begin{itemize}
    \item \textbf{Login Issues:} Check your username/password. Use \textbf{Forgot Password} if needed.
    \item \textbf{File Upload Fails:} Ensure the file is within allowed size and format.
    \item \textbf{Page Not Loading:} Refresh the page or check your internet connection.
\end{itemize}

\section{Support}
For assistance, contact the PUCON Support Team at:
\begin{itemize}
    \item Email: \href{mailto:support@pucon.example.com}{support@pucon.example.com}
    \item Phone: +63-XXX-XXX-XXXX
\end{itemize}

\appendix
\section{Appendix A: Screenshots}
\begin{center}
    % \includegraphics[width=0.8\textwidth]{assets/architecture_dia_pupsched.png}
    \captionof{figure}{Architecture diagram (optional).}
\end{center}

\section{Appendix B: Change Log}
\begin{longtable}{p{3cm} p{10cm}}
\textbf{Version} & \textbf{Changes} \\
1.0 & Initial release of the user manual. \\
\end{longtable}

\end{document}

% ----- END CLEAN VERSION -----
